% --- Archon physics formalization source begin ---
% archon:physics
% archon:covers PhyXMiniProblems/problem_phyx_mini_0639.lean
% archon:source-report reports/phyx_mini/problem_phyx_mini_0639.source.json

\chapter{Physics problem phyx\_mini\_0639}
\label{ch:PhyXMiniProblems_problem_phyx_mini_0639}

\paragraph{Problem source.}
\#\# Physical scenario

Electrons pass through the parallel electrodes shown in figure with a speed of \$5.0 \textbackslash{}times 10\textasciicircum{}6\$ \$m/s\$. Assume that the magnetic field is confined to the region between the electrodes.

\#\# Question

What magnetic field strength will allow the electrons to pass through without being deflected?

\#\# Answer choices

- A: A: \$4.0 \textbackslash{}times 10\textasciicircum{}\{-3\} T\$

- B: B: \$5.0 \textbackslash{}times 10\textasciicircum{}\{-3\} T\$

- C: C: \$6.0 \textbackslash{}times 10\textasciicircum{}\{-3\} T\$

- D: D: \$7.0 \textbackslash{}times 10\textasciicircum{}\{-3\} T\$

\#\# Figure caption (auxiliary; use the image as primary evidence)

The image shows a diagram of a parallel-plate capacitor connected to a power source. Here are the details:

- **Power Source:** There is a symbol for a battery on the left with a voltage of 200 V.
- **Parallel Plates:** Two horizontal plates are depicted, one above the other, connected to the battery, indicating a capacitor setup.
- **Dimensions:**
  - The vertical separation between the plates is labeled as 8.0 mm.
  - The horizontal length of the plates is labeled as 4.0 cm.
- **Particle:** On the right side, there is a small circled negative sign, with a green arrow pointing left, indicating the movement of a negative charge (such as an electron).

The configuration suggests a scenario of a charged particle moving between the plates in the presence of an electric field.

\#\# Dataset metadata

Magnetism; Physical Model Grounding Reasoning, Multi-Formula Reasoning

\paragraph{Recorded answer/context.}
B: B: \$5.0 \textbackslash{}times 10\textasciicircum{}\{-3\} T\$

\paragraph{Figure/image path.}
/root/proposal\_for\_physic/science-mango/phyx\_mini\_run/phyx\_data/test\_image/639.png

\paragraph{Formalization target.}
create a compiling Lean file with sorry bodies at `PhyXMiniProblems/problem\_phyx\_mini\_0639.lean`. The Lean declarations must preserve the physical quantities, dimensions or dimensional roles, figure labels, governing-law hypotheses, and final relation expressed by this problem.
Use Mathlib/Physlib names found through LeanExplore where available. If a physics API is missing, introduce faithful local abstractions rather than scalar placeholder aliases.

\begin{theorem}[Physics formalization target]
\label{thm:physics:phyx_mini_0639:target}
\lean{PhyXMiniProblems.ProblemPhyXMini0639.problem_phyx_mini_0639}
\uses{def:physics:phyx-mini-0639:phyxminiproblems-problemphyxmini0639-magneticfieldstrengthinteslas, def:physics:phyx-mini-0639:phyxminiproblems-problemphyxmini0639-parallelplateelectronsetup, def:physics:phyx-mini-0639:phyxminiproblems-problemphyxmini0639-matchesproblemandfigurereadouts, def:physics:phyx-mini-0639:phyxminiproblems-problemphyxmini0639-hasphysicalparameters, def:physics:phyx-mini-0639:phyxminiproblems-problemphyxmini0639-magneticfieldisconfinedtoplategap, def:physics:phyx-mini-0639:phyxminiproblems-problemphyxmini0639-fieldshaveuniformmagnitudeinplategap, def:physics:phyx-mini-0639:phyxminiproblems-problemphyxmini0639-satisfiesparallelplateelectricfieldlaw, def:physics:phyx-mini-0639:phyxminiproblems-problemphyxmini0639-satisfiesundeflectedlorentzforcebalance, def:physics:phyx-mini-0639:phyxminiproblems-problemphyxmini0639-recordeddatasetanswer, def:physics:phyx-mini-0639:phyxminiproblems-problemphyxmini0639-answermatchesmagneticfield}
The assigned autoformalize agent should translate this physics problem into Lean declarations in the covered file, with theorem and lemma proof bodies written as `by sorry`.
\end{theorem}
\begin{proof}
This is an autoformalization task, not a proof task. Produce statements that can later be proved by the physics prover without weakening the physical model.
\end{proof}

% --- Archon declaration topology begin ---
\section{Declaration topology}
\begin{definition}[potential Difference Dimension]
  \label{def:physics:phyx-mini-0639:phyxminiproblems-problemphyxmini0639-potentialdifferencedimension}
  \lean{PhyXMiniProblems.ProblemPhyXMini0639.potentialDifferenceDimension}
  The dimension M L² T⁻² C⁻¹ of electric potential difference
\end{definition}
\begin{proof}
  This is the declared model definition of potential Difference Dimension.
\end{proof}
\begin{definition}[electric Field Strength Dimension]
  \label{def:physics:phyx-mini-0639:phyxminiproblems-problemphyxmini0639-electricfieldstrengthdimension}
  \lean{PhyXMiniProblems.ProblemPhyXMini0639.electricFieldStrengthDimension}
  The dimension M L T⁻² C⁻¹ of electric-field strength
\end{definition}
\begin{proof}
  This is the declared model definition of electric Field Strength Dimension.
\end{proof}
\begin{definition}[magnetic Field Strength Dimension]
  \label{def:physics:phyx-mini-0639:phyxminiproblems-problemphyxmini0639-magneticfieldstrengthdimension}
  \lean{PhyXMiniProblems.ProblemPhyXMini0639.magneticFieldStrengthDimension}
  The dimension M T⁻¹ C⁻¹ of magnetic-field strength
\end{definition}
\begin{proof}
  This is the declared model definition of magnetic Field Strength Dimension.
\end{proof}
\begin{definition}[Length Quantity]
  \label{def:physics:phyx-mini-0639:phyxminiproblems-problemphyxmini0639-lengthquantity}
  \lean{PhyXMiniProblems.ProblemPhyXMini0639.LengthQuantity}
  A nonnegative, unit-independent physical length
\end{definition}
\begin{proof}
  This is the declared model definition of Length Quantity.
\end{proof}
\begin{definition}[Speed Quantity]
  \label{def:physics:phyx-mini-0639:phyxminiproblems-problemphyxmini0639-speedquantity}
  \lean{PhyXMiniProblems.ProblemPhyXMini0639.SpeedQuantity}
  A nonnegative, unit-independent speed
\end{definition}
\begin{proof}
  This is the declared model definition of Speed Quantity.
\end{proof}
\begin{definition}[Charge Magnitude Quantity]
  \label{def:physics:phyx-mini-0639:phyxminiproblems-problemphyxmini0639-chargemagnitudequantity}
  \lean{PhyXMiniProblems.ProblemPhyXMini0639.ChargeMagnitudeQuantity}
  A nonnegative electric-charge magnitude
\end{definition}
\begin{proof}
  This is the declared model definition of Charge Magnitude Quantity.
\end{proof}
\begin{definition}[Potential Difference Quantity]
  \label{def:physics:phyx-mini-0639:phyxminiproblems-problemphyxmini0639-potentialdifferencequantity}
  \lean{PhyXMiniProblems.ProblemPhyXMini0639.PotentialDifferenceQuantity}
  \uses{def:physics:phyx-mini-0639:phyxminiproblems-problemphyxmini0639-potentialdifferencedimension}
  A nonnegative potential-difference magnitude
\end{definition}
\begin{proof}
  This is the declared model definition of Potential Difference Quantity.
\end{proof}
\begin{definition}[Electric Field Strength Quantity]
  \label{def:physics:phyx-mini-0639:phyxminiproblems-problemphyxmini0639-electricfieldstrengthquantity}
  \lean{PhyXMiniProblems.ProblemPhyXMini0639.ElectricFieldStrengthQuantity}
  \uses{def:physics:phyx-mini-0639:phyxminiproblems-problemphyxmini0639-electricfieldstrengthdimension}
  A nonnegative electric-field magnitude
\end{definition}
\begin{proof}
  This is the declared model definition of Electric Field Strength Quantity.
\end{proof}
\begin{definition}[Magnetic Field Strength Quantity]
  \label{def:physics:phyx-mini-0639:phyxminiproblems-problemphyxmini0639-magneticfieldstrengthquantity}
  \lean{PhyXMiniProblems.ProblemPhyXMini0639.MagneticFieldStrengthQuantity}
  \uses{def:physics:phyx-mini-0639:phyxminiproblems-problemphyxmini0639-magneticfieldstrengthdimension}
  A nonnegative magnetic-field magnitude
\end{definition}
\begin{proof}
  This is the declared model definition of Magnetic Field Strength Quantity.
\end{proof}
\begin{definition}[length Readout]
  \label{def:physics:phyx-mini-0639:phyxminiproblems-problemphyxmini0639-lengthreadout}
  \lean{PhyXMiniProblems.ProblemPhyXMini0639.lengthReadout}
  \uses{def:physics:phyx-mini-0639:phyxminiproblems-problemphyxmini0639-lengthquantity}
  Read a physical length in a selected length unit
\end{definition}
\begin{proof}
  This is the declared model definition of length Readout.
\end{proof}
\begin{definition}[speed Readout]
  \label{def:physics:phyx-mini-0639:phyxminiproblems-problemphyxmini0639-speedreadout}
  \lean{PhyXMiniProblems.ProblemPhyXMini0639.speedReadout}
  \uses{def:physics:phyx-mini-0639:phyxminiproblems-problemphyxmini0639-speedquantity}
  Read a speed in coherent selected length and time units
\end{definition}
\begin{proof}
  This is the declared model definition of speed Readout.
\end{proof}
\begin{definition}[charge In Coulombs]
  \label{def:physics:phyx-mini-0639:phyxminiproblems-problemphyxmini0639-chargeincoulombs}
  \lean{PhyXMiniProblems.ProblemPhyXMini0639.chargeInCoulombs}
  \uses{def:physics:phyx-mini-0639:phyxminiproblems-problemphyxmini0639-chargemagnitudequantity}
  SI readout of an electric-charge magnitude, in coulombs
\end{definition}
\begin{proof}
  This is the declared model definition of charge In Coulombs.
\end{proof}
\begin{definition}[potential Difference In Volts]
  \label{def:physics:phyx-mini-0639:phyxminiproblems-problemphyxmini0639-potentialdifferenceinvolts}
  \lean{PhyXMiniProblems.ProblemPhyXMini0639.potentialDifferenceInVolts}
  \uses{def:physics:phyx-mini-0639:phyxminiproblems-problemphyxmini0639-potentialdifferencequantity}
  SI readout of a potential difference, in volts
\end{definition}
\begin{proof}
  This is the declared model definition of potential Difference In Volts.
\end{proof}
\begin{definition}[electric Field Strength In Volts Per Meter]
  \label{def:physics:phyx-mini-0639:phyxminiproblems-problemphyxmini0639-electricfieldstrengthinvoltspermeter}
  \lean{PhyXMiniProblems.ProblemPhyXMini0639.electricFieldStrengthInVoltsPerMeter}
  \uses{def:physics:phyx-mini-0639:phyxminiproblems-problemphyxmini0639-electricfieldstrengthquantity}
  SI readout of electric-field strength, in volts per metre
\end{definition}
\begin{proof}
  This is the declared model definition of electric Field Strength In Volts Per Meter.
\end{proof}
\begin{definition}[magnetic Field Strength In Teslas]
  \label{def:physics:phyx-mini-0639:phyxminiproblems-problemphyxmini0639-magneticfieldstrengthinteslas}
  \lean{PhyXMiniProblems.ProblemPhyXMini0639.magneticFieldStrengthInTeslas}
  \uses{def:physics:phyx-mini-0639:phyxminiproblems-problemphyxmini0639-magneticfieldstrengthquantity}
  SI readout of magnetic-field strength, in teslas
\end{definition}
\begin{proof}
  This is the declared model definition of magnetic Field Strength In Teslas.
\end{proof}
\begin{definition}[magnetic Field Strength In Milliteslas]
  \label{def:physics:phyx-mini-0639:phyxminiproblems-problemphyxmini0639-magneticfieldstrengthinmilliteslas}
  \lean{PhyXMiniProblems.ProblemPhyXMini0639.magneticFieldStrengthInMilliteslas}
  \uses{def:physics:phyx-mini-0639:phyxminiproblems-problemphyxmini0639-magneticfieldstrengthquantity, def:physics:phyx-mini-0639:phyxminiproblems-problemphyxmini0639-magneticfieldstrengthinteslas}
  Read magnetic-field strength in milliteslas
\end{definition}
\begin{proof}
  This is the declared model definition of magnetic Field Strength In Milliteslas.
\end{proof}
\begin{definition}[Particle Species]
  \label{def:physics:phyx-mini-0639:phyxminiproblems-problemphyxmini0639-particlespecies}
  \lean{PhyXMiniProblems.ProblemPhyXMini0639.ParticleSpecies}
  Particle species distinguished in this problem
\end{definition}
\begin{proof}
  This is the declared model definition of Particle Species.
\end{proof}
\begin{definition}[Charge Sign]
  \label{def:physics:phyx-mini-0639:phyxminiproblems-problemphyxmini0639-chargesign}
  \lean{PhyXMiniProblems.ProblemPhyXMini0639.ChargeSign}
  Sign of a particle's electric charge
\end{definition}
\begin{proof}
  This is the declared model definition of Charge Sign.
\end{proof}
\begin{definition}[Electrode Label]
  \label{def:physics:phyx-mini-0639:phyxminiproblems-problemphyxmini0639-electrodelabel}
  \lean{PhyXMiniProblems.ProblemPhyXMini0639.ElectrodeLabel}
  The two horizontal electrodes visible in the supplied image
\end{definition}
\begin{proof}
  This is the declared model definition of Electrode Label.
\end{proof}
\begin{definition}[Electrode Polarity]
  \label{def:physics:phyx-mini-0639:phyxminiproblems-problemphyxmini0639-electrodepolarity}
  \lean{PhyXMiniProblems.ProblemPhyXMini0639.ElectrodePolarity}
  Electrical polarity shown for an electrode
\end{definition}
\begin{proof}
  This is the declared model definition of Electrode Polarity.
\end{proof}
\begin{definition}[Electrode Arrangement]
  \label{def:physics:phyx-mini-0639:phyxminiproblems-problemphyxmini0639-electrodearrangement}
  \lean{PhyXMiniProblems.ProblemPhyXMini0639.ElectrodeArrangement}
  The relative geometry of the two electrodes
\end{definition}
\begin{proof}
  This is the declared model definition of Electrode Arrangement.
\end{proof}
\begin{definition}[Spatial Direction]
  \label{def:physics:phyx-mini-0639:phyxminiproblems-problemphyxmini0639-spatialdirection}
  \lean{PhyXMiniProblems.ProblemPhyXMini0639.SpatialDirection}
  Named Cartesian directions needed to interpret the side-view figure
\end{definition}
\begin{proof}
  This is the declared model definition of Spatial Direction.
\end{proof}
\begin{definition}[Parallel Electrode Figure]
  \label{def:physics:phyx-mini-0639:phyxminiproblems-problemphyxmini0639-parallelelectrodefigure}
  \lean{PhyXMiniProblems.ProblemPhyXMini0639.ParallelElectrodeFigure}
  \uses{def:physics:phyx-mini-0639:phyxminiproblems-problemphyxmini0639-lengthquantity, def:physics:phyx-mini-0639:phyxminiproblems-problemphyxmini0639-potentialdifferencequantity, def:physics:phyx-mini-0639:phyxminiproblems-problemphyxmini0639-particlespecies, def:physics:phyx-mini-0639:phyxminiproblems-problemphyxmini0639-electrodelabel, def:physics:phyx-mini-0639:phyxminiproblems-problemphyxmini0639-electrodepolarity, def:physics:phyx-mini-0639:phyxminiproblems-problemphyxmini0639-electrodearrangement, def:physics:phyx-mini-0639:phyxminiproblems-problemphyxmini0639-spatialdirection}
  Typed content of the supplied raster figure. The three dimensionful label fields are physical labels, not bare scalar stand-ins. Their numerical unit readouts are imposed separately by MatchesProblemAndFigureReadouts
\end{definition}
\begin{proof}
  This is the declared model definition of Parallel Electrode Figure.
\end{proof}
\begin{definition}[Parallel Plate Electron Setup]
  \label{def:physics:phyx-mini-0639:phyxminiproblems-problemphyxmini0639-parallelplateelectronsetup}
  \lean{PhyXMiniProblems.ProblemPhyXMini0639.ParallelPlateElectronSetup}
  \uses{def:physics:phyx-mini-0639:phyxminiproblems-problemphyxmini0639-lengthquantity, def:physics:phyx-mini-0639:phyxminiproblems-problemphyxmini0639-speedquantity, def:physics:phyx-mini-0639:phyxminiproblems-problemphyxmini0639-chargemagnitudequantity, def:physics:phyx-mini-0639:phyxminiproblems-problemphyxmini0639-potentialdifferencequantity, def:physics:phyx-mini-0639:phyxminiproblems-problemphyxmini0639-electricfieldstrengthquantity, def:physics:phyx-mini-0639:phyxminiproblems-problemphyxmini0639-magneticfieldstrengthquantity, def:physics:phyx-mini-0639:phyxminiproblems-problemphyxmini0639-particlespecies, def:physics:phyx-mini-0639:phyxminiproblems-problemphyxmini0639-chargesign, def:physics:phyx-mini-0639:phyxminiproblems-problemphyxmini0639-spatialdirection, def:physics:phyx-mini-0639:phyxminiproblems-problemphyxmini0639-parallelelectrodefigure}
  The physical apparatus and its independent observables. In particular, magneticFieldStrength is not defined from an answer choice or from the desired 5 mT value
\end{definition}
\begin{proof}
  This is the declared model definition of Parallel Plate Electron Setup.
\end{proof}
\begin{definition}[Matches Problem And Figure Readouts]
  \label{def:physics:phyx-mini-0639:phyxminiproblems-problemphyxmini0639-matchesproblemandfigurereadouts}
  \lean{PhyXMiniProblems.ProblemPhyXMini0639.MatchesProblemAndFigureReadouts}
  \uses{def:physics:phyx-mini-0639:phyxminiproblems-problemphyxmini0639-lengthreadout, def:physics:phyx-mini-0639:phyxminiproblems-problemphyxmini0639-speedreadout, def:physics:phyx-mini-0639:phyxminiproblems-problemphyxmini0639-potentialdifferenceinvolts, def:physics:phyx-mini-0639:phyxminiproblems-problemphyxmini0639-parallelplateelectronsetup}
  Problem-statement and primary-image evidence. The two equivalent separation readouts keep both the printed 8.0 mm label and its SI calibration available to the later plate-field law. No magnetic-field value or answer label occurs in these data
\end{definition}
\begin{proof}
  This is the declared model definition of Matches Problem And Figure Readouts.
\end{proof}
\begin{definition}[Has Physical Parameters]
  \label{def:physics:phyx-mini-0639:phyxminiproblems-problemphyxmini0639-hasphysicalparameters}
  \lean{PhyXMiniProblems.ProblemPhyXMini0639.HasPhysicalParameters}
  \uses{def:physics:phyx-mini-0639:phyxminiproblems-problemphyxmini0639-lengthreadout, def:physics:phyx-mini-0639:phyxminiproblems-problemphyxmini0639-speedreadout, def:physics:phyx-mini-0639:phyxminiproblems-problemphyxmini0639-chargeincoulombs, def:physics:phyx-mini-0639:phyxminiproblems-problemphyxmini0639-potentialdifferenceinvolts, def:physics:phyx-mini-0639:phyxminiproblems-problemphyxmini0639-electricfieldstrengthinvoltspermeter, def:physics:phyx-mini-0639:phyxminiproblems-problemphyxmini0639-magneticfieldstrengthinteslas, def:physics:phyx-mini-0639:phyxminiproblems-problemphyxmini0639-parallelplateelectronsetup}
  Positivity and non-vacuity conditions for the physical experiment
\end{definition}
\begin{proof}
  This is the declared model definition of Has Physical Parameters.
\end{proof}
\begin{definition}[Field Is Confined To]
  \label{def:physics:phyx-mini-0639:phyxminiproblems-problemphyxmini0639-fieldisconfinedto}
  \lean{PhyXMiniProblems.ProblemPhyXMini0639.FieldIsConfinedTo}
  A spacetime-dependent vector field vanishes outside a specified region
\end{definition}
\begin{proof}
  This is the declared model definition of Field Is Confined To.
\end{proof}
\begin{definition}[Magnetic Field Is Confined To Plate Gap]
  \label{def:physics:phyx-mini-0639:phyxminiproblems-problemphyxmini0639-magneticfieldisconfinedtoplategap}
  \lean{PhyXMiniProblems.ProblemPhyXMini0639.MagneticFieldIsConfinedToPlateGap}
  \uses{def:physics:phyx-mini-0639:phyxminiproblems-problemphyxmini0639-parallelplateelectronsetup, def:physics:phyx-mini-0639:phyxminiproblems-problemphyxmini0639-fieldisconfinedto}
  The problem's explicit support assumption: the magnetic field vanishes away from the region between the electrodes
\end{definition}
\begin{proof}
  This is the declared model definition of Magnetic Field Is Confined To Plate Gap.
\end{proof}
\begin{definition}[Fields Have Uniform Magnitude In Plate Gap]
  \label{def:physics:phyx-mini-0639:phyxminiproblems-problemphyxmini0639-fieldshaveuniformmagnitudeinplategap}
  \lean{PhyXMiniProblems.ProblemPhyXMini0639.FieldsHaveUniformMagnitudeInPlateGap}
  \uses{def:physics:phyx-mini-0639:phyxminiproblems-problemphyxmini0639-electricfieldstrengthinvoltspermeter, def:physics:phyx-mini-0639:phyxminiproblems-problemphyxmini0639-magneticfieldstrengthinteslas, def:physics:phyx-mini-0639:phyxminiproblems-problemphyxmini0639-parallelplateelectronsetup}
  The scalar strengths stored in the setup are the uniform norms of the Physlib vector fields within the plate gap
\end{definition}
\begin{proof}
  This is the declared model definition of Fields Have Uniform Magnitude In Plate Gap.
\end{proof}
\begin{definition}[Satisfies Parallel Plate Electric Field Law]
  \label{def:physics:phyx-mini-0639:phyxminiproblems-problemphyxmini0639-satisfiesparallelplateelectricfieldlaw}
  \lean{PhyXMiniProblems.ProblemPhyXMini0639.SatisfiesParallelPlateElectricFieldLaw}
  \uses{def:physics:phyx-mini-0639:phyxminiproblems-problemphyxmini0639-lengthreadout, def:physics:phyx-mini-0639:phyxminiproblems-problemphyxmini0639-potentialdifferenceinvolts, def:physics:phyx-mini-0639:phyxminiproblems-problemphyxmini0639-electricfieldstrengthinvoltspermeter, def:physics:phyx-mini-0639:phyxminiproblems-problemphyxmini0639-parallelplateelectronsetup}
  Uniform parallel-plate electrostatics: the field points from the positive upper plate to the negative lower plate and has magnitude E = ΔV / d. This law contains no magnetic-field result
\end{definition}
\begin{proof}
  This is the declared model definition of Satisfies Parallel Plate Electric Field Law.
\end{proof}
\begin{definition}[Satisfies Undeflected Lorentz Force Balance]
  \label{def:physics:phyx-mini-0639:phyxminiproblems-problemphyxmini0639-satisfiesundeflectedlorentzforcebalance}
  \lean{PhyXMiniProblems.ProblemPhyXMini0639.SatisfiesUndeflectedLorentzForceBalance}
  \uses{def:physics:phyx-mini-0639:phyxminiproblems-problemphyxmini0639-speedreadout, def:physics:phyx-mini-0639:phyxminiproblems-problemphyxmini0639-chargeincoulombs, def:physics:phyx-mini-0639:phyxminiproblems-problemphyxmini0639-electricfieldstrengthinvoltspermeter, def:physics:phyx-mini-0639:phyxminiproblems-problemphyxmini0639-magneticfieldstrengthinteslas, def:physics:phyx-mini-0639:phyxminiproblems-problemphyxmini0639-parallelplateelectronsetup}
  For the negatively charged, left-moving electron, a magnetic field directed out of the page gives a downward magnetic force, opposite the upward electric force. An undeflected path therefore has equal force magnitudes q E = q v B. This is the general Lorentz-force balance, not the requested numerical value of B
\end{definition}
\begin{proof}
  This is the declared model definition of Satisfies Undeflected Lorentz Force Balance.
\end{proof}
\begin{definition}[Answer Choice]
  \label{def:physics:phyx-mini-0639:phyxminiproblems-problemphyxmini0639-answerchoice}
  \lean{PhyXMiniProblems.ProblemPhyXMini0639.AnswerChoice}
  Labels of the four magnetic-field choices printed in the question
\end{definition}
\begin{proof}
  This is the declared model definition of Answer Choice.
\end{proof}
\begin{definition}[answer Strength In Milliteslas]
  \label{def:physics:phyx-mini-0639:phyxminiproblems-problemphyxmini0639-answerstrengthinmilliteslas}
  \lean{PhyXMiniProblems.ProblemPhyXMini0639.answerStrengthInMilliteslas}
  \uses{def:physics:phyx-mini-0639:phyxminiproblems-problemphyxmini0639-answerchoice}
  Magnetic-field strength, in milliteslas, displayed beside each choice
\end{definition}
\begin{proof}
  This is the declared model definition of answer Strength In Milliteslas.
\end{proof}
\begin{definition}[recorded Dataset Answer]
  \label{def:physics:phyx-mini-0639:phyxminiproblems-problemphyxmini0639-recordeddatasetanswer}
  \lean{PhyXMiniProblems.ProblemPhyXMini0639.recordedDatasetAnswer}
  \uses{def:physics:phyx-mini-0639:phyxminiproblems-problemphyxmini0639-answerchoice}
  The answer label recorded in the source dataset
\end{definition}
\begin{proof}
  This is the declared model definition of recorded Dataset Answer.
\end{proof}
\begin{definition}[Answer Matches Magnetic Field]
  \label{def:physics:phyx-mini-0639:phyxminiproblems-problemphyxmini0639-answermatchesmagneticfield}
  \lean{PhyXMiniProblems.ProblemPhyXMini0639.AnswerMatchesMagneticField}
  \uses{def:physics:phyx-mini-0639:phyxminiproblems-problemphyxmini0639-magneticfieldstrengthinmilliteslas, def:physics:phyx-mini-0639:phyxminiproblems-problemphyxmini0639-parallelplateelectronsetup, def:physics:phyx-mini-0639:phyxminiproblems-problemphyxmini0639-answerchoice, def:physics:phyx-mini-0639:phyxminiproblems-problemphyxmini0639-answerstrengthinmilliteslas}
  A displayed choice matches the independently modeled physical magnetic-field strength when their millitesla readouts agree
\end{definition}
\begin{proof}
  This is the declared model definition of Answer Matches Magnetic Field.
\end{proof}
\begin{lemma}[electric Field Strength eq twenty Five Thousand volts Per Meter]
  \label{lem:physics:phyx-mini-0639:phyxminiproblems-problemphyxmini0639-electricfieldstrength-eq-twentyfivethousand-voltspermeter}
  \lean{PhyXMiniProblems.ProblemPhyXMini0639.electricFieldStrength_eq_twentyFiveThousand_voltsPerMeter}
  \uses{def:physics:phyx-mini-0639:phyxminiproblems-problemphyxmini0639-electricfieldstrengthinvoltspermeter, def:physics:phyx-mini-0639:phyxminiproblems-problemphyxmini0639-parallelplateelectronsetup, def:physics:phyx-mini-0639:phyxminiproblems-problemphyxmini0639-matchesproblemandfigurereadouts, def:physics:phyx-mini-0639:phyxminiproblems-problemphyxmini0639-satisfiesparallelplateelectricfieldlaw}
  The 200 V potential difference across 8.0 mm gives 25000 V/m
\end{lemma}
\begin{proof}
  The conclusion follows from the governing relations and figure readouts named in the statement of electric Field Strength eq twenty Five Thousand volts Per Meter.
\end{proof}
\begin{lemma}[required Magnetic Field Strength eq five Milliteslas]
  \label{lem:physics:phyx-mini-0639:phyxminiproblems-problemphyxmini0639-requiredmagneticfieldstrength-eq-fivemilliteslas}
  \lean{PhyXMiniProblems.ProblemPhyXMini0639.requiredMagneticFieldStrength_eq_fiveMilliteslas}
  \uses{def:physics:phyx-mini-0639:phyxminiproblems-problemphyxmini0639-magneticfieldstrengthinteslas, def:physics:phyx-mini-0639:phyxminiproblems-problemphyxmini0639-parallelplateelectronsetup, def:physics:phyx-mini-0639:phyxminiproblems-problemphyxmini0639-matchesproblemandfigurereadouts, def:physics:phyx-mini-0639:phyxminiproblems-problemphyxmini0639-hasphysicalparameters, def:physics:phyx-mini-0639:phyxminiproblems-problemphyxmini0639-satisfiesparallelplateelectricfieldlaw, def:physics:phyx-mini-0639:phyxminiproblems-problemphyxmini0639-satisfiesundeflectedlorentzforcebalance}
  The plate-field law and zero-force condition give B = E / v = 25000 / (5 * 10\textasciicircum{}6) = 5 * 10\textasciicircum{}-3 T
\end{lemma}
\begin{proof}
  The conclusion follows from the governing relations and figure readouts named in the statement of required Magnetic Field Strength eq five Milliteslas.
\end{proof}
% --- Archon declaration topology end ---
% --- Archon physics formalization source end ---
